\section{Policy Gradient -- REINFORCE}
\label{sec:policy-gradient}

Drugim algorytmem uczenia przez wzmacnianie, który zaimplementowałem,
jest Policy Gradient oparty na metodzie REINFORCE
\cite{sutton_barto_rl}. W przeciwieństwie do DQN model nie estymuje
wartości poszczególnych akcji. Sieć neuronowa bezpośrednio definiuje
politykę, czyli rozkład prawdopodobieństwa akcji dla otrzymanego stanu.

Podobnie jak w przypadku DQN agent otrzymuje wyłącznie
\texttt{UTHObservation}. Dzięki temu podczas podejmowania decyzji
nie ma dostępu do kart krupiera ani innych informacji ukrytych
w pełnym stanie środowiska.

\subsection{Sieć polityki i wybór akcji}

Sieć polityki przyjmuje wektor utworzony przez wybrany
\texttt{StateEncoder}. Liczba wejść wynosi więc 373 dla State A
oraz 393 dla State B. Zastosowałem dwie warstwy ukryte po
256 neuronów z funkcją aktywacji ReLU oraz warstwę wyjściową
o sześciu neuronach odpowiadających wszystkim akcjom.

Sieć zwraca surowe wartości, czyli logity. Przed utworzeniem
rozkładu prawdopodobieństwa maskuję akcje niedozwolone w aktualnej
fazie rozdania. Dla legalnych akcji prawdopodobieństwa wyznaczam
za pomocą funkcji softmax.

Podczas treningu agent losuje akcję zgodnie z otrzymanym rozkładem.
Pozwala to na naturalną eksplorację bez stosowania osobnego
mechanizmu epsilon-greedy. Podczas ewaluacji wykorzystuję natomiast
politykę deterministyczną i wybieram legalną akcję o największym
prawdopodobieństwie.

\subsection{Uczenie metodą REINFORCE}

Jedno rozdanie Ultimate Texas Hold'em traktuję jako jeden epizod.
Dla każdej decyzji zapisuję zakodowany stan, wybraną akcję,
maskę akcji legalnych oraz otrzymaną nagrodę. Po zakończeniu
epizodu wyznaczam zwrot dla każdego kroku

\begin{equation}
    G_t =
    \sum_{k=t}^{T}
    \gamma^{k-t} r_k.
\end{equation}

W eksperymentach przyjąłem $\gamma = 1$, ponieważ optymalizowanym
celem jest wynik finansowy całego rozdania, bez preferowania nagrody
otrzymanej we wcześniejszym punkcie decyzji.

Parametry polityki aktualizuję zgodnie z klasycznym estymatorem
REINFORCE. Dla batcha zawierającego $B$ kompletnych epizodów
minimalizowana funkcja straty ma postać

\begin{equation}
    L(\theta)
    =
    -\frac{1}{B}
    \sum_{i=1}^{B}
    \sum_t
    G_{i,t}
    \log
    \pi_\theta(a_{i,t} \mid s_{i,t}).
\end{equation}

Zwroty obliczam niezależnie dla każdego epizodu, a parametry sieci
pozostają niezmienione podczas zbierania całego batcha. Dopiero po
zebraniu 32 rozdań wykonuję jeden krok optymalizatora Adam.
Takie podejście ogranicza wariancję aktualizacji w porównaniu
z wykonywaniem osobnego kroku optymalizacji po każdym rozdaniu.

Domyślna konfiguracja wykorzystuje learning rate równy
$0{,}001$, batch size równy 32 oraz warstwy ukryte
$(256,256)$. Nie stosuję sieci wartości, baseline'u, regularyzacji
entropii ani normalizacji zwrotów. Pozostawiam w ten sposób
REINFORCE jako względnie prostą metodę Policy Gradient, którą
mogę bezpośrednio porównać z DQN.

\subsection{Reprodukowalność i diagnostyka}

Z jednego głównego seeda deterministycznie wyprowadzam osobne
seedy dla inicjalizacji sieci, środowiska oraz losowania akcji.
Pozwala mi to powtarzać trening bez sprzęgania źródeł losowości.

Podczas treningu zapisuję między innymi wynik epizodu, liczbę
wykonanych kroków, wartość funkcji straty i normę gradientu.
Dodatkowo dla stałego zestawu obserwacji kontrolnych rejestruję
rozkład prawdopodobieństwa akcji oraz znormalizowaną entropię
polityki. Wykorzystuję te dane do oceny przebiegu uczenia
i wykrywania sytuacji, w których polityka zbyt szybko koncentruje
się na pojedynczej akcji.

Wyuczoną sieć mogę zapisać jako checkpoint wraz z konfiguracją
treningu, wariantem reprezentacji stanu oraz rodzajem funkcji
nagrody. Dzięki temu zapisany model można później jednoznacznie
odtworzyć i wykorzystać podczas ewaluacji.